\scanpage{419}
\mspage{331}
% archive sheet number: 418

% 1–5
Ce lemme implique, comme on l'a vu à la page 312, l'implication $b\Longrightarrow d$ dans le diagramme d'implications \unsure{(*)|($\alpha$)|(a)} page 310, lequel est donc à présent prouvé.

% 6–7
Pour achever de prouver le théorème 4.13.15, il reste à prouver les implications
% 8–9
\diagram{\begin{tikzcd}[ampersand replacement=\&,row sep=0.5em]
\deleted{\unsure{d}}\unsure{c'|d'}\arrow[dr,Rightarrow] \& \&\\
\& \deleted{\unsure{b}}\unsure{a'|b'}\arrow[r,Rightarrow] \& a)\\
b\arrow[ur,Rightarrow] \& \&
\end{tikzcd}}

% 10
Démontrons donc $[\unsure{b)|b'}\ \text{ou}\ \unsure{c')|d')}]\Longrightarrow a')$.

% 10–12, interline at right
\inserted{Il reste à prouver que $\underline S$ est $\unsure{\pi|\varpi}$-accessible dans $M$, i.e.} \deleted{Notons} \deleted{Comme $\underline S$ est déjà supposée str. pleine dans $M$, il nous reste à prouver} \deleted{que \unsure{$\underline S$}, \unsure{le} \illegible{1}}
% 12–15
que le foncteur pl. fidèle \deleted{$\pi$ \unsure{étant} \illegible{1} \unsure{pour} $M$,}
% 14
\begingroup\csname tagsleft@true\endcsname
\[
\operatorname{Ind}_\pi(\underline S(\pi))\longrightarrow\deleted{\unsure{\underline S}}\underline M
\tag*{(*)}
\]
\endgroup
% 15
a pour image essentielle \deleted{$\unsure{S}$} $\underline S$.

% 15–20
\deleted{$M_\pi/x$ est grande devant $\pi$, donc \inserted{la sous-catégorie str. pleine} $\underline S(\pi)/x$ est cofinale, elle est elle-$=$ grande devant $\pi$, a fortiori filtrante.

Donc c) implique que si $\unsure{x^-|x}\in\operatorname{Ob}\underline S$, alors on a $x=\varinjlim_{i\in\underline S(\pi)/x}x_i$, \deleted{ce qui implique} $\underline S(\pi)/x$ gde devant $\pi$. Comme}

% 21–22
Donc il faut prouver que si $x\in\operatorname{Ob} M$, alors
% 23–26
\begingroup\csname tagsleft@true\endcsname
\[
\begin{aligned}
 x\in\operatorname{Ob}\underline S\quad&\overset{\unsure{(?)|(!)}}{\Longleftrightarrow}\quad x\ \text{est dans l'image essentielle}\\
 &\hspace{4em}\text{de }\operatorname{Ind}_\pi(\underline S(\pi)),\ \text{i.e. }x\\
 &\hspace{4em}\text{est isomorphe à }\varinjlim_{i\in I}x_i,\\
 &\hspace{4em}\text{les }x_i\ \text{dans }\operatorname{Ob}(\underline S(\pi)\unsure{\,|)} ,\ I\ \text{grd}\\
 &\hspace{4em}\text{devant }\pi.
\end{aligned}
\tag*{$\substack{(*)\\\deleted{\unsure{(*)}}}$}
\]
\endgroup
